\section*{Transferencia del momento y de la exponencial}
Se mantienen las paradas de v3/06. Sea \(X=Qe^{-L}\), \(g_0=\log X(B_0)\).
Sobre historias detenidas la densidad de soluciones frente a la ley original
es \(W=\exp(A-g_0)X\). Vale cero en las historias no completables.
En \(H_j\), \(X=\exp(g_T+\Gamma+g^\flat)\), \(|g_T|\le C/j\),
\(\sup Y_D^+\le C\) por v3/06 y \(|\Gamma|\le C/j\) por la cota de par y el
recuento determinista de la parte 02 anterior. La constante de la cota superior
en \(H_j\) es de clase: el corte determinista usa \(j\) como tamaño y los márgenes
se fijaron antes que \(b\). Además \(|A-g_0|=O_{b,\delta}(1/n)\).
Luego
\[
 \mathbb E_\mu[|g^\flat|;H_j]\le C\mathbb E_\nu[|g^\flat|;H_j].
\]
Para la exponencial que entra con signo en la identidad se tiene exactamente
\[
 X\,|e^{-g^\flat}-1|=e^{g_T+\Gamma}|1-e^{g^\flat}|.
\]
Como \(|e^y-1|\le e^{y^+}|y|\), su media bajo \(\mu\) tiene la misma cota.
Para el signo positivo, en \(H_j\) también basta la cota de \(Y_D^+\).

Sea \(F_b\) la llegada regular completable y \(T=F_b\setminus H_j\).
Es una subclase de la llegada larga regular. V3/06 da
\(\mathbb E_\nu[1+X;T]\le C e^{-2j}\).
En esa clase \(|g_T|+|\Gamma|\le C/j\) por la pertenencia de la máscara a la
clase de torres (parte 04) y las cotas de las configuraciones individuales.
El Corolario 6.8 del artículo de la etapa final da
\((\log X)^+\le C_0+\beta j\) en la clase exterior.
Para \(0<X\le1\) se usa \(X|\log X|\le1/e\), y para \(X>1\) se usa
\(X|\log X|\le X(C_0+\beta j)\). Por tanto
\[
 \mathbb E_\mu[|g^\flat|+|e^{-g^\flat}-1|;T]\le C(1+j)e^{-2j}.
\]
En la exponencial negativa se ha usado
\(X|e^{-g^\flat}-1|\le e^{g_T+\Gamma}+X\).
Si se solicita la exponencial positiva en \(T\), el mismo cálculo da
\(O(e^{-(2-\beta)j})\) en la clase de carga con \(\beta<2\).
La conclusión necesaria para la identidad, independiente de esta última variante, es
\[
 \limsup_{n\to\infty}n\mathbb E_\mu
 \bigl[|g^\flat|+|e^{-g^\flat}-1|;F_b\bigr]\le C_\delta b .
\]
Se ha probado una media integrada al orden de límites pedido.
La desigualdad puntual E3\(\flat\) en todas las máscaras sigue sin decidirse.
